% Part: first-order-logic 
% Chapter: axiomatic-deduction 
% Section: rules-and-proofs

\documentclass[../../../include/open-logic-section]{subfiles}

\begin{document}

\iftag{FOL}
      {\olfileid{fol}{axd}{rul}}
      {\olfileid{pl}{axd}{rul}}

\olsection{القواعد و\usetoken{P}{derivation}}

\begin{explain}
  قد تكون أنساق ال!!{derivation}s البديهية أبسط أنساق ال!!{derivation}
  في المنطق. ف!!^a{derivation} ليس إلا متتالية من ال!!{formula}s.
  ولكي تُعد المتتالية !!a{derivation}، يجب أن تكون كل !!{formula} فيها
  إما حالة من بديهية، وإما أن تلزم بقاعدة استدلال عن !!{formula}s واحدة
  أو أكثر تسبقها في المتتالية. وتقوم !!^a{derivation} بــ!!{derive}s
  آخر !!{formula} فيها.
\end{explain}

\begin{defn}[!!^{derivability}]
إذا كانت $\Gamma$ مجموعة من !!{formula}s اللغة $\Lang L$، فإن
\emph{!!{derivation}} من~$\Gamma$ هو متتالية منتهية $!A_1$،
\dots،~$!A_n$ من ال!!{formula}s، يتحقق لكل $i \le n$ واحد مما يأتي:
\begin{enumerate}
\item $!A_i \in \Gamma$؛ أو
\item $!A_i$ بديهية؛ أو
\item تلزم $!A_i$ عن $!A_j$ ما (وعن $!A_k$) حيث $j < i$ (و
  $k < i$) بقاعدة استدلال.
\end{enumerate}
\end{defn}

يعتمد ما يُعد !!{derivation} صحيحًا على قواعد الاستدلال التي نجيزها (وطبعًا
على ما نتخذه بديهيات). وقاعدة الاستدلال عبارة شرطية تخبرنا بأنه في ظروف
معينة تكون خطوة~$!A_i$ في !!a{derivation} خطوة استدلال صحيحة.

\begin{defn}[قاعدة الاستدلال]
تعطي \emph{قاعدة الاستدلال} شرطًا كافيًا لما يُعد خطوة استدلال صحيحة في
!!a{derivation} من~$\Gamma$.
\end{defn}

فمثلًا، بما أن كل متتالية أحادية العنصر $!A$ تحقق $!A \in \Gamma$ تُعد
بوضوح !!a{derivation}، فقد تكون العبارة الآتية قاعدة استدلال بسيطة جدًا:
\begin{quote}
  إذا كان $!A \in \Gamma$، فإن $!A$ دائمًا خطوة استدلال صحيحة في كل
  !!{derivation} من~$\Gamma$.
\end{quote}
وبالمثل، إذا كانت $!A$ إحدى البديهيات، فإن $!A$ وحدها !!a{derivation}،
ولذلك تكون العبارة الآتية أيضًا قاعدة استدلال:
\begin{quote}
  إذا كانت $!A$ بديهية، فإن $!A$ خطوة استدلال صحيحة.
\end{quote}
ويصبح الأمر أكثر إثارة إذا أحالت قاعدة الاستدلال إلى !!{formula}s تظهر
قبل الخطوة موضع النظر. وتسمى القاعدة الآتية \emph{قاعدة إثبات المقدَّم}:
\begin{quote}
  إذا ظهرت $!B \lif !A$ و$!B$ في موضع أعلى من ال!!{derivation}،
  فإن~$!A$ خطوة استدلال صحيحة.
\end{quote}
إذا كانت هذه قاعدة الاستدلال الوحيدة، كان تعريفنا لل!!{derivation} أعلاه
مكافئًا لما يأتي: تكون $!A_1$، \dots،~$!A_n$ !!a{derivation} إذا وفقط
إذا تحقق لكل $i \le n$ واحد مما يأتي:
\begin{enumerate}
\item $!A_i \in \Gamma$؛ أو
\item $!A_i$ بديهية؛ أو
\item لبعض $j < i$ تكون $!A_j$ هي $!B \lif !A_i$، ولبعض $k < i$ تكون
  $!A_k$ هي~$!B$.
\end{enumerate}
يقول الشرط الأخير إن $!A_i$ تلزم عن~$!A_j$ ($!B \lif !A_i$) و$!A_k$
($!B$) بقاعدة إثبات المقدَّم. فإذا استطعنا الانتقال من $1$ إلى~$n$،
وفي كل مرة نجد !!a{formula}~$!A_i$ إما تنتمي إلى~$\Gamma$، وإما بديهية،
وإما تخبرنا قاعدة استدلال بأنها خطوة استدلال صحيحة، فإن المتتالية كلها تُعد
!!{derivation} صحيحًا.

\begin{defn}[!!^{derivability}]
تكون ال!!{formula}~$!A$ \emph{!!{derivable}} من $\Gamma$، ونكتب
$\Gamma \Proves !A$، إذا وُجد !!a{derivation} من $\Gamma$ ينتهي
بــ$!A$.
\end{defn}

\begin{defn}[المبرهنات]
تكون !!^a{formula}~$!A$ \emph{مبرهنة} إذا وُجد !!a{derivation}
لــ~$!A$ من المجموعة الخالية. ونكتب $\Proves !A$ إذا كانت $!A$ مبرهنة،
و$\Proves/ !A$ إذا لم تكن كذلك.
\end{defn}


\end{document}
